\chapter{Evaluation Methods}
\label{chap:evaluation}
To evaluate a model, we look beyond the loss function. 
There are a few reasons for this, and that largest is simply that most evaluation methods are \textit{not} differentiable.
Therefore our most useful and expressive evaluation metrics (a suite of error representations) cannot be used during training.
We can use them afterwards though.

\section{Classification Evaluation}
\label{sec:evaluation-classification-evaluation-the-confusion-matrix}
For binary classification, predictions fall into true positive (TP), false positive (FP), true negative (TN) and false negative (FN).
\subsection{Confusion Matrix}
\begin{center}
    \small
    \begin{tabular}{@{}lcc@{}}
    \toprule
     & Actual positive & Actual negative \\
    \midrule
    Predicted positive & TP & FP \\
    Predicted negative & FN & TN \\
    \bottomrule
    \end{tabular}
\end{center}
Accuracy $(\mathrm{TP}+\mathrm{TN})/N_{\text{examples}}$ can hide rare-class failure because the majority class dominates the count.

\subsection{Precision, Recall and F1}
\label{sec:evaluation-precision-recall-and-f1}
\begin{align}
    \underbrace{P}_{\text{precision: correctness of alerts}}
    &=\frac{\mathrm{TP}}{\mathrm{TP}+\mathrm{FP}},\\
    \underbrace{R}_{\text{recall: coverage of positives}}
    &=\frac{\mathrm{TP}}{\mathrm{TP}+\mathrm{FN}},\\
    F_\beta
    &=(1+\beta^2)\frac{PR}{\beta^2P+R}.
\end{align}
Larger $\beta$ weights recall more strongly and smaller positive $\beta$ weights precision more strongly.
F1 ignores TN directly, so it is not automatically suitable when correctly rejecting negatives is itself important.

\begin{calloutspecialcase}
    The $F_1$ score is $F_\beta$ with $\beta=1$.
\end{calloutspecialcase}


\subsection{ROC curves, AUC-ROC and Precision--Recall}
\label{sec:evaluation-roc-curves-auc-roc-and-precision-recall}
The Receiver Operating Characteristic (ROC) curves compares true-positive rate against false-positive rate at different thresholds.
ROC-Area-Under-Curve (AUC) is the probability that a random positive is ranked above a random negative.

\figpdf[.9\textwidth]{roc_precision_recall}{ROC and precision--recall curves expose different threshold trade-offs.}

For rare events, a small false-positive rate can still create many more false alerts than true alerts.
Precision--Recall curves expose that directly, so rare-event systems should not rely on accuracy or ROC-AUC alone.

\begin{calloutpractical}
    Rare positives: report precision, recall and a precision--recall curve, then choose threshold from false-alarm versus miss cost.
\end{calloutpractical}


\section{Folded Validation}
\label{sec:evaluation-cross-validation-and-fold-choice}
As ML practitioners, we like to get statistical answers to our questions, and yet this model evaluation section trains one model and evaluates once with multiple metrics.
This is \textit{not} a strong statistical methodology, and so we train repeats on subsets of the data. 
These subsets are folds, and each fold is validated once while the others train the model.

The mean fold score is
\[
    \bar s
    =
    \frac{1}{N_{\text{folds}}}
    \sum_{i=1}^{N_{\text{folds}}}
    \underbrace{s_i}_{\text{score on held-out fold }i}.
\]
\begin{calloutpractical}
    Fold scores are correlated because their training sets overlap.
\end{calloutpractical}

\begin{figure}[htbp]
    \centering
    \begin{tikzpicture}[x=.97cm,y=.54cm]
    \foreach \r in {0,...,4}{
    \node[anchor=east,font=\small] at (-.16,-\r) {Run \the\numexpr\r+1\relax};
    \foreach \c in {0,...,4}{
    \ifnum\c=\r \fill[amber!85] (\c*1.7,-\r-.35) rectangle ++(1.5,.75);
    \else \fill[deepblue!82] (\c*1.7,-\r-.35) rectangle ++(1.5,.75);\fi
    \node[white,font=\small] at (\c*1.7+.75,-\r+.02) {\the\numexpr\c+1\relax};}}
    \fill[deepblue!82] (1,-5.7) rectangle ++(.35,.3);
    \node[anchor=west,font=\scriptsize] at (1.45,-5.55) {Training};
    \fill[amber!85] (4.2,-5.7) rectangle ++(.35,.3);
    \node[anchor=west,font=\scriptsize] at (4.65,-5.55) {Validation};
    \end{tikzpicture}
    \caption{Five-fold cross-validation: each fold is held out once.}
\end{figure}
    
\begin{tabularx}{\linewidth}{@{}L{.24\linewidth}YY@{}}
    \toprule
    Split strategy & Problem it Solves & Main leakage controlled \\
    \midrule
    Random K-fold & Examples are approximately independent and similarly distributed & Dependence on one arbitrary split \\
    Stratified K-fold & Class proportions vary strongly & Unrepresentative class balance \\
    Group K-fold & Repeated measurements share subject/run/instrument identity & Same group appearing on both sides \\
    Forward / time-series split & Future prediction & Training on future observations \\
    Nested CV & Hyperparameter selection with limited data & Reusing selection folds for final assessment \\
    \bottomrule
\end{tabularx}

\begin{calloutpractical}
Split by the independent unit: keep repeated measurements from one subject/run together, and preserve time order for forecasting.
\end{calloutpractical}


\section{Ranking metrics}
\label{sec:evaluation-retrieval-and-ranking-metrics}

\paragraph{Recall@K.}
If $N_{\text{relevant}}$ items exist and $N_{\text{relevant in top K}}$ survive the first stage,
\[
    \underbrace{\operatorname{Recall@K}}_{\text{candidate coverage}}
    =
    \frac{N_{\text{relevant in top K}}}{N_{\text{relevant}}}.
\]
A downstream reranker cannot recover an item that retrieval never produced.

\paragraph{MRR.}
Mean Reciprocal Rank rewards putting the first relevant result near the top:
\[
    \operatorname{MRR}
    =
    \frac{1}{N_{\text{queries}}}
    \sum_q
    \frac{1}{\underbrace{\operatorname{rank}_{q,\text{first relevant}}}_{\text{where useful information first appears}}}.
\]

\paragraph{NDCG.}
Discounted Cumulative Gain gives top positions more weight and allows graded relevance:
\[
    \underbrace{\operatorname{NDCG@K}}_{\text{normalised ranking quality}}
    =
    \frac{
    \sum_{r=1}^{K}
    \frac{2^{\operatorname{relevance}_r}-1}{\log_2(r+1)}
    }{
    \underbrace{\operatorname{IDCG@K}}_{\text{best possible ordering for this query}}
    }.
\]
The logarithmic discount is why swapping ranks 1 and 2 matters much more than swapping ranks 101 and 102.
\figpdf[.76\textwidth]{ndcg_discount}{NDCG's position discount falls quickly near the top, so high-ranking mistakes receive most of the metric weight.}

\begin{calloutpractical}
Evaluate candidate generation with coverage-oriented metrics such as Recall@K, then evaluate the later ordering with NDCG/MRR; one scalar for the whole funnel can hide which stage failed.
\end{calloutpractical}

\section{Generative Evaluation}
\label{sec:evaluation-generative-and-rag-evaluation}
\subsection{Golden Dataset}
A \textbf{golden dataset} acts like a test suite. 
A perfect, human derived set of examples allows reruns whenever the system changes, in the style of unit tests in software engineering.
Because models have biases, \textbf{LLM as a Judge} is limited, and requires calibration, usually from a golden set.

\subsection{Rag Evaluation}
\textbf{RAGAS}-style evaluation focuses on context precision/recall, answer relevancy and faithfulness.
\textbf{HELM}-style benchmarking instead asks broader base-model capability questions across standardised scenarios.

\begin{calloutmemoryhook}
    For RAG, separate ``did we retrieve the evidence?'' from ``did the generator use the evidence correctly?''.
\end{calloutmemoryhook}

\section{A/B Testing (Online Evaluation)}
\label{sec:evaluation-online-experiments-and-ab-testing}
How do you know whether a model is actually better in practice than an alternative?
You compare them in the real-world setting. 
We need a statistically strong signal to justify rolling out an iteration, and so we use statistics to answer the question.

We turn the deployment into an experiment, testing the delta:
\[
    \underbrace{\widehat{\Delta}}_{\text{estimated treatment effect}}
    =
    \underbrace{\bar y_{\text{treatment}}}_{\text{new system}}
    -
    \underbrace{\bar y_{\text{control}}}_{\text{old system}}.
\]
\textbf{Persistent assignment}, such as a stable hash of user identity, keeps one user in one condition and avoids mixing treatment exposure.

The null hypothesis is commonly $H_0:\Delta=0$.
A p-value asks how extreme the observed statistic would be under the null.
The standard-error scaling in \cref{sec:foundations-confidence-intervals} provides a sample-size approximation called the \textbf{Minimum Detectable Effect} (MDE):
\[
    \underbrace{\operatorname{MDE}}_{\text{smallest reliably detectable effect}}
    \propto
    \frac{1}{\sqrt{N_{\text{independent cases}}}}
    \quad\text{(all else equal)}.
\]
Halving the target MDE therefore costs roughly four times as many independent observations under the same noise/decision assumptions.

\begin{center}
    \small
    \begin{tabularx}{.9\linewidth}{@{}L{.25\linewidth}YY@{}}
    \toprule
     & Test says no effect & Test says effect \\
    \midrule
    No real effect & Correct & Type I error \\
    Real effect & Type II error & Correct detection \\
    \bottomrule
    \end{tabularx}
\end{center}

Repeatedly checking an ordinary fixed-horizon p-value and stopping the first time it crosses a threshold changes the testing procedure and inflates false-positive risk.
Use a predeclared horizon or a sequential method whose stopping rule is part of the design.

Ranking systems can also use \textbf{interleaving}: mix candidates from two rankers into one list so the same user/context provides a sensitive relative comparison.
Interleaving is useful for direct ranker preference but does not replace product-level A/B metrics such as retention or revenue.

\begin{calloutdegreeofexperiment}
    Vary assignment unit, traffic split, metric, MDE, test duration, stopping rule, segmentation and whether competing ranking systems are compared by A/B assignment or interleaving.
\end{calloutdegreeofexperiment}
